\documentclass[10pt,a4paper]{amsart}
\usepackage[margin=19mm]{geometry}
\usepackage{amsmath,amssymb,mathtools}
\usepackage[scr=rsfs,cal=euler]{mathalfa}
\usepackage{xfrac}
\usepackage[unicode,hidelinks,bookmarks=false]{hyperref}
\newcommand{\ie}{i.e.\ }
\newcommand{\cf}{cf.\ }
\newcommand{\resp}{respectively\ }
\newcommand{\Subsection}{\subsection}
\providecommand{\nobreakdash}{\nobreak\hbox{-}}
\setlength{\emergencystretch}{2em}
\multlinegap0pt
\usepackage{amssymb}
\newtheorem{thm}{Theorem}
\newcommand{\cF}{\mathcal{F}}
\title{Shape Theory $\&$ TDA via the Atiyah--Molino Reconstruction}
\author{N. C. Combe, H. K. Nencka}
\date{Source: arXiv:2608.15696v1 (CC BY 4.0)}
\begin{document}
\maketitle

\noindent\emph{Source and licence.} Shape Theory $\&$ TDA via the Atiyah--Molino Reconstruction. Version arXiv:2608.15696v1 (CC BY 4.0), distributed by arXiv under the Creative Commons Attribution 4.0 International licence, http://creativecommons.org/licenses/by/4.0/, as recorded in the arXiv licence metadata for this exact version and shown on the abstract page https://arxiv.org/abs/2608.15696v1. Full licence notice: \texttt{licenses/arxiv-2608.15696-CC-BY-4.0.txt}.\par\smallskip

\noindent\textit{Editorial scope.} Counted unit: the article's thm thm (source0.tex lines 184--201). The article's complete original proof of it is delivered with it (source0.tex lines 203--222, one \texttt{proof} environment of 20 lines). No mathematics is written, repaired or re-derived: the statement and the proof are the article's own lines, in the article's own notation, and no retained line is substituted or re-flowed.  Retained uncounted as the source-local context the proof needs: lines 102--102; lines 160--160.  Nothing was omitted from any retained span, so the package carries no omission record.  No retained line contains a citation, so the wrapper carries no bibliography and no bibliography entry.  No retained line contains a figure environment, an image-inclusion command or drawing code, so no image is delivered and the package declares no asset dependency.  Editorial formatting only, in addition: an AMS \texttt{amsart} preamble that restates the article's own declarations and macros used by the retained lines, namely the environments \texttt{thm} on separate counters, together with the standard packages those lines need.  The retained environments print in this document as Theorem 1 in order of appearance, whereas the article prints the same items with the numbering of its own full document.  The complete native source of the article as distributed in the arXiv e-print for this exact version is bundled unchanged as \texttt{sources/207772/source0.tex}.
\section{Recollections on Vaisman--Neifeld's Setting}\label{S:VN}

\subsubsection{}\label{S:QNI} Let \( M \) be the foliated reconstruction space described in Sec. \ref{S:VN}, equipped with:

\begin{thm}
Suppose that \( M \) is defined as above (Sec. \ref{S:QNI})  and equipped with transverse foliations \( \mathcal{F}_1, \mathcal{F}_2 \), Haantjes tensors \( H_1, H_2 \) and the algebra $(\mathcal{A} , \star)$. Then,
\begin{enumerate}
    \item The algebra \( \mathcal{A} \) is {\bf associative} if and only if the Haantjes tensors \( H_1, H_2 \) vanish identically. Equivalently:
    \[
    \mathcal{A} \text{ is associative} \iff H_i = 0 \quad (i = 1, 2).
    \]
    \item If \( H_i \neq 0 \), the algebraic structure of \( \mathcal{A} \) is governed by a Moufang-like identity:
    \[
    (a \star b) \star (c \star a) = a \star (b \star c) \star a, \quad \forall a,b,c \in \mathcal{A}
    \]
    reflecting the curvature of \( \nabla_1 \otimes \nabla_2 \).
    \item \begin{itemize} 
    \item If the {\bf associative case} holds, then the reconstruction problem has a unique solution \( O \in M \), and \( M \) is globally biholomorphic to \( \mathbb{C}P^3 \). 
    \item If the {\bf non-associative case} holds, then solutions form a quasigroupoid under \( \star \), with non-unique reconstructions parameterized by the cohomology class \( [H_i] \in H^1(M, TM) \).
    \end{itemize}
\end{enumerate}
\end{thm}

\begin{proof}~
\begin{itemize}
\item[(1)] \textbf{Haantjes Tensor as an Obstruction.} The vanishing of the Haantjes tensor $H_i$	
ensures that the corresponding distributions $D_i$	
are integrable, thereby endowing the foliation 
$F_i$ with a Lie foliation structure. In this case, the induced connection $\nabla_i$	
 is necessarily flat, and the Leibniz rule forces the underlying algebra 
$\mathcal{A}$ to be associative. Reciprocally, if $\mathcal{A}$ is associative, the holonomy correction must be path-independent. Non-vanishing $H_i$ (torsion) introduces path-dependence in parallel transport via the torsion-twisted Lie bracket, causing non-associativity in $\star$. Thus the associativity forces $H_i=0$.
Therefore, \[\mathcal{A} \text{ is associative} \iff H_i = 0 \quad (i = 1, 2).\]
\item[]
\item[(2)] \textbf{Curvature and the Emergence of Non-Associativity.} The vanishing of $H_i$ ensures associativity by eliminating path-dependence. Conversely, when \(H_i \neq 0\), the curvature \(F_{\nabla_i}\) measures deviation from associativity. Indeed, the curvature  \(F_{\nabla_i}\) encodes the local holonomy and parallel transport deviations. For the tensor product connection  \(\nabla=\nabla_1 \otimes \nabla_2\), the curvature \(F_{\nabla}\) governs the failure of associativity in $\mathcal{A}$. Specifically, the associator $(a,b,c)=(a\star b)\star c-a\star (b\star c) $ is non-zero. The Bianchi identity for the tensor product connection \(\nabla_1 \otimes \nabla_2\), is a differential constraint on curvature: \[d^{\nabla}F_{\nabla}+[F_{\nabla}\wedge F_{\nabla}]=0,\] where $d^{\nabla}$ is the covariant exterior derivative. Under the imposed involution symmetry (an anti-automorphism $a\mapsto \overline{a}$ satisfying $ \overline{a\star b}= \overline{b}\star \overline{a}$), the Bianchi identity forces the associator to satisfy a symmetrized relation. Naturally, this symmetry reduces the general non-associativity to the Moufang identity: \[(a\star b)\star (c\star a)=a\star (b\star c)\star a, \] which is a weaker, symmetric form of associativity, thereby reflecting the inherent non-associative structure.
\item[]
\item[(3)] \begin{itemize}
\item \textbf{Global Implications.} In the associative setting, the triviality of holonomy guarantees a unique global parallel transport, thereby facilitating a coherent reconstruction. In contrast, non-associativity introduces nontrivial monodromy, which obstructs global uniqueness. In such circumstances, the quasigroupoid structure emerges from the Ehresmann connection defined on the product foliation \(\mathcal{F}_1 \times \mathcal{F}_2\).

\item \textbf{Quasigroupoid}.  The non-unique reconstructions (solutions) correspond to morphisms between objects, which are leaves of the foliations $\cF_1, \cF_2$.
The operation $\star$  is defined only for solutions sharing compatible projection data that is overlapping domains/codomains in the foliated space. The cohomology class $[H_i] \in H^1(M, TM)$ encodes deformations of the foliations, which act as automorphisms on the quasigroupoid structure. To summarise, the objects are leaves of $\cF_1, \cF_2$, morphisms correspond to reconstructions $O$, with compatible projections. The $\star$ operations correspond to composition of deformations along paths in $M$.
\end{itemize}
\end{itemize}
\end{proof}

\end{document}
